\documentclass[12pt]{article}
\usepackage[utf8]{inputenc}
\usepackage[T1]{fontenc}
\usepackage{hyperref}
\title{Yala Language Resources for K-12 Curriculum Development}
\author{Compiled by Claude Sonnet 5.5 (claude-sonnet-5-5) for Andr\'as Kornai}
\date{}
\begin{document}
\maketitle

\section{Introduction}
Yala (ISO 639-3 yba; Glottocode yala1263) is an Idomoid language of the Ogoja area of Nigeria. Blench (2019), as cited by Wikipedia, lists Ikom, Obubra and Ogoja dialects. Only Bible translation material turned up in the search results.

\section{Primary Educational Resources}
No primers, textbooks or school programmes were found.

\section{Linguistic and Reference Materials}
\begin{itemize}
\item Yala language (Wikipedia overview). \url{https://en.wikipedia.org/wiki/Yala_language}
\end{itemize}
No grammar, dictionary or orthography description was found.

\section{Religious and Community Resources}
\begin{itemize}
\item Lutheran Bible Translators, Yala Bible Translation programme. It reports a printed New Testament with Psalms, a revision of the New Testament in progress, and Old Testament drafting under way. \url{https://lbt.org/program/yala-bible-translation/}
\item Lutheran Bible Translators, Yala Language Development and Bible Translation (page and PDF). \url{https://lbt.org/?p=5329} and \url{https://lbt.org/wp-content/uploads/2020/09/Yala-Language-Development-and-BT.pdf}
\item Joshua Project profile of Yala in Nigeria, a missions profile that is secondary as a source. \url{https://joshuaproject.net/people_groups/15961}
\end{itemize}

\section{Digital and Audio Resources}
None were found.

\section{Curriculum Development Context}
The printed Yala New Testament is the only substantial text found, and it implies a working orthography that the search did not describe. Dialect variation (Ikom, Obubra, Ogoja) would need a decision on which variety to use. Classroom materials, a dictionary and a documented orthography were not found.

\end{document}